\section*{Chapitre \ref{ch:b1:rock-forming-minerals} --- Les minéraux des roches}

\begin{solution}{exo:b1:rock-forming-minerals:1}
\ce{SiO4^{4-}}~: \omterm{def:b1:rock-forming-minerals:silicate-classes}{nésosilicate}, $s = 0$. \ce{Si2O7^{6-}}~: \omterm{def:b1:rock-forming-minerals:silicate-classes}{sorosilicate},
$s = 1$. \ce{Si6O18^{12-}}~: \omterm{def:b1:rock-forming-minerals:silicate-classes}{cyclosilicate}, $s = 2$. \ce{Si2O5^{2-}}~:
\omterm{def:b1:rock-forming-minerals:silicate-classes}{phyllosilicate}, $s = 3$. \ce{SiO2}~: \omterm{def:b1:rock-forming-minerals:silicate-classes}{tectosilicate}, $s = 4$. \ce{SiO3^{2-}}
a $4 - s/2 = 3$, donc $s = 2$~: un anneau (\omterm{def:b1:rock-forming-minerals:silicate-classes}{cyclosilicate}) ou une chaîne simple
(\omterm{def:b1:rock-forming-minerals:silicate-classes}{inosilicate})~; le rapport seul ne permet pas de les distinguer.
\end{solution}

\begin{solution}{exo:b1:rock-forming-minerals:2}
Forstérite~: $2\times24.305 + 28.085 + 4\times15.999 =
\qty{140.691}{g.mol^{-1}}$~; fayalite~: $2\times55.845 + 28.085 +
4\times15.999 = \qty{203.771}{g.mol^{-1}}$. Volumes molaires~:
$140.691/3.275 = \qty{42.96}{cm^3.mol^{-1}}$ et $203.771/4.392 =
\qty{46.40}{cm^3.mol^{-1}}$.
\end{solution}

\begin{solution}{exo:b1:rock-forming-minerals:3}
$(0.78 - 0.72)/0.72 = +\qty{8.3}{\percent}$ pour \ce{Fe^{2+}}~;
$-\qty{4.2}{\percent}$ pour \ce{Ni^{2+}}~; $+\qty{15.3}{\percent}$ pour
\ce{Mn^{2+}}~; $+\qty{38.9}{\percent}$ pour \ce{Ca^{2+}}. Le fer et le nickel, à
moins de \qty{15}{\percent}, remplacent librement le magnésium~; le manganèse,
juste au-dessus, est limité~; le calcium, au-delà de \qty{30}{\percent}, n'entre
guère. Les quatre ont la même charge que le magnésium.
\end{solution}

\begin{solution}{exo:b1:rock-forming-minerals:4}
$(3.51 - 2.16)/2.16 = \qty{62.5}{\percent}$ du graphite au diamant~;
$(2.93 - 2.71)/2.71 = \qty{8.1}{\percent}$ de la calcite à l'aragonite. Les
\omterm{def:b1:rock-forming-minerals:polymorph}{polymorphes} les plus denses, le diamant et l'aragonite, sont les formes de haute
pression.
\end{solution}

\begin{solution}{exo:b1:rock-forming-minerals:5}
Avec $x = 0.30$ dans $\mathrm{Na}_{1-x}\mathrm{Ca}_x\mathrm{Al}_{1+x}\mathrm{Si}_{3-x}\mathrm{O}_8$~:
$\mathrm{(Na_{0.7}Ca_{0.3})(Al_{1.3}Si_{2.7})O_8}$. Charges~:
$0.7 + 2\times0.3 + 3\times1.3 + 4\times2.7 = 0.7 + 0.6 + 3.9 + 10.8 = +16$,
et $8\times(-2) = -16$.
\end{solution}

\begin{solution}{exo:b1:rock-forming-minerals:6}
Un groupe répété comprend deux tétraèdres partageant trois sommets (ceux dont
l'apex joint les deux chaînes) et deux en partageant deux. Les quatre tétraèdres
ont $4\times4 = 16$ sommets, dont $3 + 3 + 2 + 2 = 10$ sont partagés, soit
$10/2 = 5$ oxygènes partagés~; les $16 - 10 = 6$ autres ne sont pas partagés.
Total $5 + 6 = 11$ oxygènes pour 4 siliciums~: \ce{Si4O11}, de charge
$4\times4 - 11\times2 = -6$. La proposition donne la même chose avec
$s = 5/2$~: $4 - 5/4 = 11/4$ oxygènes par silicium, charge $5/2 - 4 = -3/2$ par
silicium.
\end{solution}

\begin{solution}{exo:b1:rock-forming-minerals:7}
Moles dans \qty{100}{g}~: \ce{SiO2} $40.20/60.083 = 0.66907$~; \ce{MgO}
$47.10/40.304 = 1.16862$~; \ce{FeO} $12.20/71.844 = 0.16981$~; \ce{NiO}
$0.40/74.692 = 0.00536$. Oxygènes~: $2\times0.66907 + 1.16862 + 0.16981 +
0.00536 = 2.68194$~; $f = 4/2.68194 = 1.49146$. Pour quatre oxygènes~: Si
\num{0.998}, Mg \num{1.743}, Fe \num{0.253}, Ni \num{0.008}. Les sites
octaédriques contiennent $1.743 + 0.253 + 0.008 = 2.004$, le site tétraédrique
\num{0.998}~: une bonne analyse, \ce{(Mg_{1.743}Fe_{0.253}Ni_{0.008})SiO4}.
$\mathrm{Fo} = 100\times1.743/(1.743 + 0.253) = 87.3$. Le facteur $f$
multiplie Mg et Fe de la même façon, donc il s'élimine dans le rapport~:
$\mathrm{Fo} = 100\,(47.10/40.304)/(47.10/40.304 + 12.20/71.844) =
100\times1.16862/1.33843 = 87.3$.
\end{solution}

\begin{solution}{exo:b1:rock-forming-minerals:8}
Positif~: $1 + 2\times3 + 3\times4 + 3 = +22$~; négatif~:
$10\times(-2) + 2\times(-1) = -22$. Un feuillet de composition \ce{Si4O10}
porterait $-4$ (\cref{prop:b1:rock-forming-minerals:si-o}, $s = 3$)~; le
remplacement d'un silicium sur quatre par de l'aluminium ajoute une charge
négative, donc \ce{(Si3Al)O10} porte $-5$. Les deux aluminiums octaédriques avec
les deux hydroxyles, $6 - 2 = +4$, en équilibrent quatre sur cinq~; l'ion
potassium entre les feuillets équilibre la cinquième.
\end{solution}

\begin{solution}{exo:b1:rock-forming-minerals:9}
(a) $M = (140.691 + 203.771)/2 = \qty{172.231}{g.mol^{-1}}$, $V = (42.96 +
46.40)/2 = \qty{44.68}{cm^3.mol^{-1}}$, $\rho = 172.231/44.68 =
\qty{3.855}{g.cm^{-3}}$~; la moyenne des masses volumiques est $(3.275 + 4.392)/2 =
\qty{3.834}{g.cm^{-3}}$, soit \qty{0.021}{g.cm^{-3}} de moins. (b) Résoudre
$203.771 + x(140.691 - 203.771) = 3.600\,\bigl(46.40 + x(42.96 - 46.40)\bigr)$~:
$203.771 - 63.080x = 167.040 - 12.384x$, donc $x = 36.731/50.696 = 0.7245$~:
l'\omterm{def:b1:rock-forming-minerals:groups}{olivine} est $\text{Fo}_{72}$ (au pour cent près).
\end{solution}

\begin{solution}{exo:b1:rock-forming-minerals:10}
Avec $M_i = \rho_i V_i$,
\[
  \rho(x) - \bigl(x\rho_1 + (1 - x)\rho_2\bigr)
  = \frac{x\rho_1V_1 + (1 - x)\rho_2V_2 - \bigl(x\rho_1 + (1 - x)\rho_2\bigr)\bigl(xV_1 + (1 - x)V_2\bigr)}{V(x)} .
\]
En développant le produit, le numérateur vaut $x\rho_1V_1 - x^2\rho_1V_1 +
(1 - x)\rho_2V_2 - (1 - x)^2\rho_2V_2 - x(1 - x)(\rho_1V_2 + \rho_2V_1)$, et
$x - x^2 = x(1 - x)$, $(1 - x) - (1 - x)^2 = x(1 - x)$, donc il vaut
$x(1 - x)(\rho_1V_1 + \rho_2V_2 - \rho_1V_2 - \rho_2V_1) =
x(1 - x)(\rho_1 - \rho_2)(V_1 - V_2)$. Pour l'\omterm{def:b1:rock-forming-minerals:groups}{olivine} (1 = forstérite), les deux
différences sont négatives, $3.275 - 4.392 = -1.117$ et $42.96 - 46.40 =
-3.44$, donc leur produit est positif et la masse volumique est au-dessus de la
droite. En $\text{Fo}_{50}$~: $0.25\times1.117\times3.44/44.68 =
\qty{0.0215}{g.cm^{-3}}$, ce qui concorde avec l'\cref{exo:b1:rock-forming-minerals:9}
aux arrondis près des volumes molaires.
\end{solution}

\begin{solution}{exo:b1:rock-forming-minerals:11}
Total \qty{99.60}{\percent}. Moles d'oxyde~: \ce{SiO2} $0.90042$,
\ce{Al2O3} $0.28344$, \ce{CaO} $0.19973$, \ce{Na2O} $0.08229$,
\ce{K2O} $0.00318$. Oxygènes~: $1.80084 + 0.85033 + 0.19973 + 0.08229 +
0.00318 = 2.93636$~; $f = 8/2.93636 = 2.72446$. Cations (deux fois les moles pour
\ce{Al2O3}, \ce{Na2O}, \ce{K2O})~: Si \num{2.453}, Al \num{1.544}, Ca
\num{0.544}, Na \num{0.448}, K \num{0.017}. Sites tétraédriques~:
$2.453 + 1.544 = 3.997$~; grand site~: $0.544 + 0.448 + 0.017 = 1.009$. Alors
$\mathrm{Al} - 1 = 0.544$, égal à Ca aux arrondis près~: chaque calcium est entré
avec un aluminium. $\mathrm{An} = 100\times0.544/1.009 = 53.9$~: une
labradorite, $\text{An}_{54}$.
\end{solution}

\begin{solution}{exo:b1:rock-forming-minerals:12}
(a) $M(\ce{Fe3O4}) = 3\times55.845 + 4\times15.999 = \qty{231.531}{g.mol^{-1}}$.
Ses trois fers, rapportés comme \ce{FeO}, pèsent $3\times71.844 = \qty{215.532}{g}$,
soit $215.532/231.531 = \qty{93.09}{\percent}$~; l'oxygène lié au
\ce{Fe^{3+}}, au-delà de ce que compte \ce{FeO}, manque au total.
(b) $93.09/71.844 = 1.29572$ mole de \ce{FeO} et d'oxygène~; $f =
4/1.29572 = 3.08708$~; quatre fers pour quatre oxygènes, \ce{Fe4O4}. La magnétite
a trois cations pour quatre oxygènes. Les oxygènes ont été sous-comptés~: chaque
\ce{Fe^{3+}} en tient un et demi, non un comme le suppose \ce{FeO}. (c) Sur trois
cations, Fe $= 3$ avec $a$ \ce{Fe^{2+}} et $b$ \ce{Fe^{3+}}~: $a + b = 3$ et
$2a + 3b = 8$, donc $b = 2$, $a = 1$~: $\mathrm{Fe^{2+}Fe^{3+}_2O_4}$. (d) Un
\ce{FeO} et un \ce{Fe2O3} par unité formulaire~: $71.844/231.531 = \qty{31.03}{\percent}$
de \ce{FeO} et $159.687/231.531 = \qty{68.97}{\percent}$ de \ce{Fe2O3}, total
\qty{100.00}{\percent}.
\end{solution}

\begin{solution}{pb:b1:rock-forming-minerals:1}
\textbf{1.} \ce{SiO2} $28.085 + 2\times15.999 = 60.083$~; \ce{FeO} 71.844~;
\ce{MnO} 70.937~; \ce{MgO} 40.304~; \ce{NiO} 74.692~; \ce{CaO} 56.077
(\si{g.mol^{-1}}).

\textbf{2.} $39.50 + 16.53 + 0.23 + 43.32 + 0.20 + 0.22 =
\qty{100.00}{\percent}$. Un total de \qty{97}{\percent} aurait suggéré un
constituant non analysé (l'eau, ou l'oxygène lié à \ce{Fe^{3+}} rapporté comme
\ce{FeO}), ou une analyse médiocre.

\textbf{3.} \ce{SiO2} $39.50/60.083 = 0.65742$~; \ce{FeO} 0.23008~; \ce{MnO}
0.00324~; \ce{MgO} 1.07483~; \ce{NiO} 0.00268~; \ce{CaO} 0.00392.

\textbf{4.} $2\times0.65742 = 1.31485$ pour \ce{SiO2}, les mêmes nombres qu'à la
question 3 pour les autres~; somme $1.31485 + 0.23008 + 0.00324 + 1.07483 +
0.00268 + 0.00392 = 2.62960$.

\textbf{5.} Chaque oxyde a un cation~: les nombres de la question 3.

\textbf{6.} $f = 4/2.62960 = 1.52114$.

\textbf{7.} Si $0.65742\times1.52114 = 1.000$~; Fe 0.350~; Mn 0.005~; Mg 1.635~;
Ni 0.004~; Ca 0.006.

\textbf{8.} $1.635 + 0.350 + 0.005 + 0.004 + 0.006 = 2.000$ pour les sites
octaédriques et 1.000 pour le site tétraédrique~: exactement les 2 et 1 idéaux.

\textbf{9.} \ce{(Mg_{1.635}Fe_{0.350}Mn_{0.005}Ni_{0.004}Ca_{0.006})SiO4}.

\textbf{10.} $2\times2.000 + 4\times1.000 = +8.000$ contre $4\times(-2) =
-8$. Chaque cation a été compté avec la charge de son oxyde, si bien que les
charges positives valent le double des oxygènes par construction~: la
vérification ne peut pas échouer. Les totaux des sites de la question 8, si.

\textbf{11.} Le nickel diffère du magnésium de
$(0.69 - 0.72)/0.72 = -\qty{4}{\percent}$, bien en deçà de \qty{15}{\percent}~:
il se substitue librement. Le calcium est $(1.00 - 0.72)/0.72 = \qty{39}{\percent}$
plus gros, au-delà de \qty{30}{\percent}~: il n'entre guère.

\textbf{12.} Du fer compté comme \ce{FeO} alors qu'une partie est du \ce{Fe^{3+}}
donne trop peu d'oxygènes par cation, donc plus de trois cations pour quatre
oxygènes (\cref{exo:b1:rock-forming-minerals:12}). La somme est
$1.000 + 2.000 = 3.000$~: pas de \ce{Fe^{3+}} significatif.

\textbf{13.} $\mathrm{Fo} = 100\times1.635/(1.635 + 0.350) = 163.5/1.985 =
82.4$.

\textbf{14.} $\mathrm{Mg\#} = 82.4$~: tout le fer est du \ce{Fe^{2+}}
(question 12), donc $\mathrm{Fe^{2+}} = \mathrm{Fe}$ et les deux définitions
coïncident.

\textbf{15.} $\text{Fo}_{82.4}\text{Fa}_{17.6}$, $x = 0.824$~; le fer remplit
\qty{17.6}{\percent} des sites octaédriques dans le modèle binaire
($0.350/2.000 = \qty{17.5}{\percent}$ avec les éléments mineurs).

\textbf{16.} Le facteur multiplie Mg et Fe de la même façon et s'élimine~:
$\mathrm{Fo} = 100\times1.07483/(1.07483 + 0.23008) = 82.4$.

\textbf{17.} $40.5/40.304 = 1.00486$ et $20.0/71.844 = 0.27838$~;
$\mathrm{Fo} = 100\times1.00486/1.28324 = 78.3$. L'\omterm{def:b1:rock-forming-minerals:groups}{olivine} tardive est plus riche
en fer.

\textbf{18.} Masses molaires~: $M = 140.691$ et \qty{203.771}{g.mol^{-1}}~;
$V = 140.691/3.275 = 42.96$ et $203.771/4.392 =
\qty{46.40}{cm^3.mol^{-1}}$.

\textbf{19.} $M = 0.824\times140.691 + 0.176\times203.771 = 115.929 +
35.864 = \qty{151.793}{g.mol^{-1}}$~; $V = 0.824\times42.96 +
0.176\times46.40 = 35.399 + 8.166 = \qty{43.565}{cm^3.mol^{-1}}$.

\textbf{20.} $\rho = 151.793/43.565 = \qty{3.484}{g.cm^{-3}}$.

\textbf{21.} $0.824\times3.275 + 0.176\times4.392 = 2.699 + 0.773 =
\qty{3.472}{g.cm^{-3}}$, soit \qty{0.012}{g.cm^{-3}} (\qty{0.3}{\percent}) au-dessous
de la valeur de mélange idéal~: proche, comme le prévoit
l'\cref{exo:b1:rock-forming-minerals:10}.

\textbf{22.} $0.12\times\qty{3484}{kg.m^{-3}}\times\qty{1}{m^3} =
\qty{418}{kg}$.

\textbf{23.} $418/2950 = \qty{14.2}{\percent}$~: l'\omterm{def:b1:rock-forming-minerals:groups}{olivine} dense représente une
plus grande part de la masse que du volume.

\textbf{24.} L'\omterm{def:b1:rock-forming-minerals:groups}{olivine} est $\text{Fo}_{82}$~: sa teneur en forstérite est de
\textbf{82\,\%}.
\end{solution}
